\documentclass[12pt,letterpaper]{article}

% ── Packages ─────────────────────────────────────────────────────────────────
\usepackage{amsmath,amssymb,amsthm}
\usepackage{mathtools}
\usepackage{geometry}
\usepackage{hyperref}
\usepackage{booktabs}
\usepackage{array}
\usepackage{xcolor}
\usepackage{framed}
\usepackage{microtype}
\usepackage[numbers,sort&compress]{natbib}

\geometry{margin=1in}

% ── Theorem environments ──────────────────────────────────────────────────────
\theoremstyle{plain}
\newtheorem{theorem}{Theorem}[section]
\newtheorem{lemma}[theorem]{Lemma}
\newtheorem{proposition}[theorem]{Proposition}
\theoremstyle{definition}
\newtheorem{definition}[theorem]{Definition}
\newtheorem{hypothesis}[theorem]{Hypothesis}
\theoremstyle{remark}
\newtheorem{remark}[theorem]{Remark}

% ── Framed boxes for key results ─────────────────────────────────────────────
\newenvironment{keyresult}{%
  \begin{framed}\noindent\textbf{Key Result.}\quad}{%
  \end{framed}}

% ── Custom commands ───────────────────────────────────────────────────────────
\newcommand{\tr}{\operatorname{Tr}}
\newcommand{\kB}{k_{\mathrm{B}}}
\newcommand{\lnTwo}{\ln 2}
\newcommand{\lP}{\ell_{P}}
\newcommand{\kt}{\tilde{\kappa}}
\newcommand{\kteff}{\tilde{\kappa}_{\mathrm{eff}}}
\newcommand{\Sent}{S_{\mathrm{ent}}}
\newcommand{\Ment}{M_{\mathrm{ent}}}
\newcommand{\asc}{\alpha_{\mathrm{screen}}}
\newcommand{\Deltaa}{\Delta a}

% ── Metadata ─────────────────────────────────────────────────────────────────
\title{\textbf{Gravitational Coupling to Entanglement Entropy Density}\\[4pt]
\large A Falsifiable Extension of Thermodynamic Gravity to Quantum-Coherent Systems}

\author{Kevin Monette\\[2pt]
\small Independent Researcher (AI-assisted research)\\
\small \texttt{kevin.monette@research.org}}

\date{February 9, 2026 \quad (repository version March 2026)}

% ─────────────────────────────────────────────────────────────────────────────
\begin{document}
\maketitle

\begin{abstract}
We derive a dimensionally consistent coupling between entanglement entropy density and
spacetime curvature, extending Jacobson's thermodynamic formulation of general
relativity~\cite{Jacobson1995} to macroscopic quantum-coherent systems.  The modified Einstein
equation takes the form
\begin{equation*}
G_{\mu\nu}
= \frac{8\pi G}{c^{4}}\,T_{\mu\nu}
+ \frac{\kt}{k_{\mathrm{B}}\lnTwo}\,\Sent\,g_{\mu\nu},
\end{equation*}
where $\Sent$ is the entanglement entropy density (bit/m$^{3}$) and $\kt$ is a dimensionless
coupling constant.  Under the \emph{Planck-scale holographic regulator hypothesis} --- the
assumption that laboratory-scale entanglement entropy couples to geometry through the same
Planck-length regulation that governs horizon entropy --- we obtain the estimate $\kt = -1/4$.
This is an \emph{estimate under an explicit hypothesis}, not a derivation from first principles.
Realistic systems exhibit environmental suppression $\kteff = -(1/4)\asc$, where $\asc$ is an
open parameter requiring derivation from open-quantum-system dynamics.

Under spherical symmetry, the modified Poisson equation gives an entanglement-induced
acceleration correction
\begin{equation*}
  \Deltaa(R) = -\frac{\kt\,c^{4}}{8\pi\kB\lnTwo}\cdot
               \frac{\Ment(<R)}{R^{2}},
\end{equation*}
where $\Ment(<R)$ is the enclosed entanglement entropy.  For bounded sources this
scales as $1/R^{2}$; a $1/R$ law requires a spatially extended source profile with
$\rho_{\mathrm{ent}}\propto 1/r^{2}$.

The signal magnitude is currently indeterminate because $\asc$ and the Planck-scale
renormalization of $\kt$ are unresolved.  The falsification criterion, however, is
independent of signal magnitude: if macroscopic quantum-coherent systems ($\geq 10^{6}$
entangled qubits) exhibit no anomalous stress-energy at pressure sensitivity
$\Delta p < 10^{-6}$\,Pa after $\geq 1000$ runs across three or more independent platforms,
then $|\kt| < 10^{-15}$, ruling out laboratory-scale relevance.

\medskip\noindent
\textbf{Ontological constraints.} Classical spacetime, standard quantum mechanics, no new
particles or modified geometry --- only a modified stress-energy source via entanglement
entropy.
\end{abstract}

\tableofcontents
\bigskip

% ─────────────────────────────────────────────────────────────────────────────
\section{Introduction: Thermodynamic Gravity and Its Extension}
\label{sec:intro}

\subsection{Jacobson's Result}

In 1995, Jacobson showed that Einstein's field equations follow from the Clausius relation
$\delta Q = T\,dS$ applied to local Rindler horizons~\cite{Jacobson1995}.  For an observer
accelerating with proper acceleration $a$, the local horizon temperature is the Unruh
temperature
\begin{equation}
  T_{U} = \frac{\hbar a}{2\pi c\kB},
  \label{eq:unruh}
\end{equation}
and the entropy associated with a horizon area element $dA$ is the Bekenstein--Hawking entropy
\begin{equation}
  dS_{BH} = \frac{\kB c^{3}}{4G\hbar}\,dA.
  \label{eq:bh-entropy}
\end{equation}
Applying $\delta Q = T_{U}\,dS_{BH}$ and identifying the energy flux with the
stress-energy tensor recovers $G_{\mu\nu} = 8\pi G c^{-4} T_{\mu\nu}$ exactly.

Gravity is therefore not fundamental in this view: it is the thermodynamic equation of state
of spacetime when entropy scales with horizon area.

\subsection{The Natural Extension}

If gravity emerges from area-law entanglement entropy at horizons, the question arises: does
volume-law entanglement entropy in a laboratory quantum system contribute to the gravitational
source term?  The horizon is absent, but the entanglement is physically real.

We propose this as a falsifiable physical hypothesis.  Supporting evidence includes
Verlinde's 2017 demonstration that entropic gravity extends beyond the strict horizon
case~\cite{Verlinde2017}, and Bose et al.'s observation of gravity-mediated entanglement
between massive quantum objects without causal horizons~\cite{Bose2017}.

% ─────────────────────────────────────────────────────────────────────────────
\section{The Modified Einstein Equation: Dimensional Rigor}
\label{sec:theory}

\subsection{Bit-to-Entropy Conversion}

``Bit'' is a counting unit, dimensionless like ``dozen.''  Physical entropy requires
conversion:
\begin{equation}
  S_{\mathrm{physical}} = I_{\mathrm{bits}}\cdot\kB\lnTwo,
  \qquad
  \kB\lnTwo = 9.570\times10^{-24}\;\text{J/K}.
  \label{eq:conversion}
\end{equation}

We define:
\begin{align}
  \Sent &\equiv \text{entanglement entropy density (bit/m}^{3}\text{),} \\
  \mathcal{S}_{\mathrm{ent}} &\equiv \Sent\cdot\kB\lnTwo \quad
    \text{[J/(K$\cdot$m$^{3}$)]:\ physical entropy density.}
\end{align}

\subsection{The Modified Field Equation}
\label{ssec:field-eq}

\begin{equation}
  \boxed{
  G_{\mu\nu}
  = \frac{8\pi G}{c^{4}}
    \!\left(
      T_{\mu\nu}
      + \kt\,\frac{c^{4}}{8\pi G\kB\lnTwo}\,\Sent\,g_{\mu\nu}
    \right)}
  \label{eq:modified-einstein}
\end{equation}

\paragraph{Dimensional verification.}
The entanglement source term must match $[T_{\mu\nu}] = \text{kg\,m}^{-1}\text{s}^{-2}$:
\begin{equation}
  \left[\kt\frac{c^{4}}{8\pi G\kB\lnTwo}\Sent\right]
  = \frac{\text{m}^{4}\text{s}^{-4}}
         {\text{m}^{3}\text{kg}^{-1}\text{s}^{-2}
          \cdot\text{kg\,m}^{2}\text{s}^{-2}\text{K}^{-1}}
    \cdot\text{K\,m}^{-3}
  = \text{kg\,m}^{-1}\text{s}^{-2}. \checkmark
  \label{eq:dim-check}
\end{equation}

\paragraph{Repulsive gravity condition.}
For a perfect fluid with energy density $\rho_{m}$ and pressure $p$:
\begin{equation}
  \rho_{\mathrm{grav}} + \frac{3p_{\mathrm{grav}}}{c^{2}}
  = \rho_{m} + \frac{3p}{c^{2}}
    + \frac{3\kt c^{2}}{8\pi G\kB\lnTwo}\,\Sent.
  \label{eq:source}
\end{equation}
For $\kt < 0$ and $\Sent > 0$, the entanglement term contributes negative effective pressure,
producing repulsive curvature without exotic matter.

% ─────────────────────────────────────────────────────────────────────────────
\section{The Extrapolation Boundary}
\label{sec:boundary}

This section is the intellectual center of the paper.

\subsection{Where Jacobson's Derivation Applies}

Jacobson's result holds exactly under three conditions:
\begin{enumerate}
  \item A causal horizon exists, defining a well-posed Unruh temperature~\eqref{eq:unruh}.
  \item Entanglement entropy scales with horizon \emph{area} (Bekenstein--Hawking law).
  \item The surface gravity $\kappa_{s}$ is well defined and uniform on the horizon.
\end{enumerate}
None of these conditions holds in a laboratory quantum system: there is no causal horizon,
no Unruh temperature, and GHZ-state entanglement entropy scales with atom \emph{number}
(volume law), not surface area.

\subsection{The Hypothesis}

\begin{hypothesis}[Holographic Regulator Extension]
\label{hyp:extension}
The coupling between entanglement entropy density and spacetime curvature, which follows
rigorously from area-law entropy at causal horizons, extends to systems with volume-law
entanglement entropy in the absence of causal horizons, regulated by the same Planck-scale
geometric structure that governs horizon area entropy.
\end{hypothesis}

Hypothesis~\ref{hyp:extension} is motivated by holographic universality and the Bose et al.\
result~\cite{Bose2017}, but it is not mathematically derivable from established physics.  Its
scientific value lies entirely in its experimental falsifiability (Section~\ref{sec:falsification}).

% ─────────────────────────────────────────────────────────────────────────────
\section{Derivation of $\kt$}
\label{sec:kappa}

\emph{All results in this section are conditional on Hypothesis~\ref{hyp:extension}.}

\subsection{Modified Clausius Relation}

Under Hypothesis~\ref{hyp:extension}, the entanglement entropy contributes an additional term
to the Clausius relation.  For an entanglement entropy density $\Sent$ (bit/m$^{3}$), the
additional entropy associated with horizon area element $dA$ is:
\begin{equation}
  dS_{\mathrm{ent}}
  = \frac{\Sent}{\kB}\cdot\frac{dV}{4\lP},
  \qquad
  dV = \lP\,dA,
  \label{eq:dSent}
\end{equation}
where $\lP = \sqrt{\hbar G/c^{3}} = 1.616\times10^{-35}$\,m is the Planck length.

\begin{remark}
Equation~\eqref{eq:dSent} uses $dV = \lP\,dA$, treating one Planck length of depth behind
the horizon as the effective thickness of the entanglement contribution.  This regulator is
borrowed from holographic principles and has no rigorous derivation for non-horizon systems.
It is the central unproven assumption of this framework.
\end{remark}

\subsection{Algebra}

Substituting $T_{U} = \hbar a/(2\pi c\kB)$ and $dV = \lP\,dA$:
\begin{equation}
  \delta Q_{\mathrm{eff}}
  = \delta Q_{BH}
    + \frac{\hbar a}{2\pi c\kB}\cdot\frac{\Sent}{\kB}\cdot\frac{dA}{4}.
  \label{eq:dQeff}
\end{equation}
Following Jacobson's identification of $\delta Q_{\mathrm{eff}} = T^{\mathrm{eff}}_{\mu\nu}
k^{\mu}d\Sigma^{\nu}$, using surface gravity $a = c^{2}\kappa_{s}$, and converting
$\Sent \to \Sent\cdot\kB\lnTwo$:
\begin{equation}
  T^{\mathrm{ent}}_{\mu\nu}
  = -\frac{c^{4}}{32\pi G}\,\Sent\,g_{\mu\nu}.
  \label{eq:Tent}
\end{equation}
Comparing with Eq.~\eqref{eq:modified-einstein}:

\begin{keyresult}
Under Hypothesis~\ref{hyp:extension} (Planck-scale holographic regulator),
\begin{equation}
  \kt = -\frac{1}{4}.
  \label{eq:kappa-result}
\end{equation}
This is an estimate motivated by thermodynamic analogy.  It is not a derivation from first
principles.
\end{keyresult}

\subsection{Environmental Screening}

Real systems are continuously decohered by their environments.  The effective coupling is:
\begin{equation}
  \kteff = -\frac{1}{4}\,\asc,
  \label{eq:kappa-eff}
\end{equation}
where $\asc \in [0,1]$ is the fraction of the ideal coupling surviving environmental
decoherence.

\begin{remark}[Open Problem 3]
The screening factor $\asc$ must be derived from the Lindblad master equation for the
specific experimental system.  The heuristic range $[10^{-4}, 10^{-2}]$ in earlier drafts
applied to few-qubit mesoscopic systems.  For macroscopic ensembles ($N \sim 10^{6}$ atoms,
cloud radius $\sim 0.1$\,mm, coherence length $\sim 1\,\mu$m), Gaussian spatial screening
yields $\asc \sim 10^{-43}$, making the macroscopic atom-interferometry experiment
ill-motivated as the primary near-term target.  Until $\asc$ is derived for the relevant
regime, the predicted signal magnitude is indeterminate.
\end{remark}

% ─────────────────────────────────────────────────────────────────────────────
\section{Radial Scaling: Resolution of OP4}
\label{sec:radial}

In the Newtonian (weak-field) limit, Eq.~\eqref{eq:modified-einstein} becomes a modified
Poisson equation for the entanglement potential $\Psi$:
\begin{equation}
  \nabla^{2}\Psi = \kappa_{\mathrm{eff}}\,\rho_{\mathrm{ent}}(R),
  \qquad
  \kappa_{\mathrm{eff}} \equiv \frac{\kt c^{4}}{2\kB\lnTwo}.
  \label{eq:poisson}
\end{equation}
Under spherical symmetry, the radial equation~\eqref{eq:poisson} integrates to:
\begin{equation}
  \Deltaa(R)
  \;\equiv\;
  -\frac{d\Psi}{dR}
  \;=\;
  \frac{\kappa_{\mathrm{eff}}}{R^{2}}\,\Ment(<R),
  \label{eq:delta-a}
\end{equation}
where the enclosed entropic source is
\begin{equation}
  \Ment(<R) = 4\pi\int_{0}^{R}\rho_{\mathrm{ent}}(r)\,r^{2}\,dr.
  \label{eq:Ment}
\end{equation}

\paragraph{Bounded sources.}
For a compact atomic ensemble (radius $R_{0}$), $\Ment(<R) \to \Ment^{\mathrm{tot}}$ for
$R > R_{0}$, giving:
\begin{equation}
  \Deltaa(R) \propto \frac{1}{R^{2}}, \qquad R > R_{0}.
  \label{eq:1/R2}
\end{equation}

\paragraph{Extended sources.}
A $1/R$ correction requires $\rho_{\mathrm{ent}}(r) \propto 1/r^{2}$ so that
$\Ment(<R) \propto R$, giving $\Deltaa(R) \propto 1/R$.  This is \emph{not} a universal
consequence of the modified Poisson equation; it encodes a specific density profile assumption.

\begin{remark}
The $1/R$ formula in earlier source documents treated this scaling as universal.
Equation~\eqref{eq:delta-a} resolves the discrepancy: the $1/R$ behavior encodes a hidden
assumption that $\Ment(<R) \propto R$.  For a compact atomic ensemble, the correct exterior
prediction is $1/R^{2}$.  \textbf{OP4 is closed.}
\end{remark}

% ─────────────────────────────────────────────────────────────────────────────
\section{Current Experimental Bounds}
\label{sec:bounds}

No experiment has been specifically designed to test this coupling.  Table~\ref{tab:bounds}
summarizes upper bounds from null results in related experiments.  The MICROSCOPE bound is
omitted pending derivation of the explicit mapping from satellite observables to $\kt$
(see Section~\ref{ssec:microscope}).

\begin{table}[h]
\centering
\caption{Experimental upper bounds on $|\kt|$ from null results.}
\label{tab:bounds}
\begin{tabular}{@{}lll@{}}
\toprule
Experiment & Constraint & Reference \\
\midrule
Gravity-mediated entanglement & $|\kt| < 3\times10^{-9}$ & Bose et al.~\cite{Bose2017}$^{\dagger}$ \\
Atom interferometry (dual-species) & Constrains $\eta$ at $10^{-12}$ & Asenbaum et al.~\cite{Asenbaum2020} \\
\bottomrule
\multicolumn{3}{l}{$^{\dagger}$\footnotesize Mapping from gravity-entanglement bound to $\kt$ bound to be verified.} \\
\end{tabular}
\end{table}

\subsection{MICROSCOPE and the Missing Mapping}
\label{ssec:microscope}

MICROSCOPE constrains the E\"{o}tv\"{o}s parameter at the $10^{-15}$ level~\cite{MICROSCOPE2022}.
Both test masses are classical solids with $\Sent \approx 0$: the entanglement source term
contributes essentially equally (zero) to both trajectories, yielding $\eta \approx 0$
independently of $\kt$.  \emph{A MICROSCOPE-derived bound on $\kt$ requires an explicit
argument mapping microscopic entanglement structure to satellite-observable gravitational
anomalies.  Until that argument is constructed, no MICROSCOPE bound on $\kt$ is stated.}

% ─────────────────────────────────────────────────────────────────────────────
\section{Revised Experimental Protocol: NISQ-Scale Test}
\label{sec:experiment}

\subsection{Motivation for the Revised Target}

The original Section~5 proposed $N \geq 10^{6}$ $^{87}$Rb atoms in a GHZ-like state.  Given
$\asc \sim 10^{-43}$ for a macroscopic BEC (Section~\ref{sec:kappa}), this target is not
well-motivated as the primary near-term experiment.  We revise the proposal to a 50--100 qubit
NISQ-scale platform.

\subsection{Platform Specifications}

\begin{table}[ht]
\centering\small
\caption{Minimum NISQ platform requirements for the revised experiment.}
\label{tab:nisq}
\begin{tabular}{@{}lp{3.5cm}p{4cm}@{}}
\toprule
Parameter & Requirement & Current capability (2025) \\
\midrule
Qubit count & 50--100 & $>100$ (multiple platforms) \\
Entanglement depth & $\geq 10$ qubits (verified) & Achievable with 2-qubit gates \\
Gate fidelity (2Q) & $\geq 99\%$ & 99--99.5\% (best platforms) \\
Coherence time & $\geq 100\,\mu$s & $100\,\mu$s--s \\
State verification & Witnesses or tomography & Standard \\
\bottomrule
\end{tabular}
\end{table}

\subsection{Protocol}

\paragraph{Branch~A (entangled).}
Prepare an $N$-qubit state with verified entanglement depth $k \geq 10$: GHZ-like, cluster
state, or spin-squeezed state.  Verify entanglement using witnesses and subsystem purity.

\paragraph{Branch~B (reference).}
Apply controlled dephasing (random single-qubit $Z$ rotations) to destroy inter-qubit
entanglement while preserving qubit count, spatial geometry, and gross energy budget.
Verify entanglement depth $\leq 1$ (product state).

\paragraph{Observable.}
Measure a differential observable (phase shift, frequency shift, or effective acceleration
proxy) that could be sensitive to an entanglement-dependent coupling.  Systematically vary
entanglement depth between runs and look for any reproducible dependence.

\subsection{Signal Magnitude}

The predicted signal magnitude is indeterminate.  The structure of the signal is given by
Eq.~\eqref{eq:delta-a}; any numeric forecast requires first resolving $\asc$ (Open
Problem~3) and the Planck-scale renormalization of $\kt$ (Open Problem~10).  The experiment
is designed not to claim detectability but to produce a bound on
$|\kteff| \cdot \Ment^{\mathrm{tot}} / R^{2}$ --- a meaningful experimental result regardless
of outcome.

\subsection{Systematic Error Budget}

\begin{table}[ht]
\centering\small
\caption{Dominant systematics in the NISQ-scale protocol and their mitigations.}
\label{tab:systematics}
\begin{tabular}{@{}lp{3cm}p{4.5cm}@{}}
\toprule
Systematic & Origin & Mitigation \\
\midrule
Imperfect state prep. & Gate errors, init.\ fidelity & Entanglement verification per run \\
Decoherence drift & Time-dependent noise & Alternating A-B-A-B; calibration \\
Geometry coupling & Device layout & Model $\rho_{\mathrm{ent}}(x)$ per architecture \\
Platform cross-talk & Qubit--qubit coupling & Symmetrized pulses; null circuits \\
Ambiguous screening & $\asc$ unknown & State signal as indeterminate \\
\bottomrule
\end{tabular}
\end{table}

\subsection{External Benchmarks}

The 2020 Asenbaum et al.\ experiment~\cite{Asenbaum2020} used a dual-species
$^{85}$Rb/$^{87}$Rb interferometer with $T \sim 2$\,s free fall:
\begin{equation}
  \eta = [1.6 \pm 1.8\,(\text{stat}) \pm 3.4\,(\text{syst})] \times 10^{-12},
  \label{eq:eta}
\end{equation}
with per-shot resolution $\sim\!1.4\times10^{-11}\,g$.  Any entanglement-induced effect
behaving as a composition-dependent fifth force must respect this bound.

% ─────────────────────────────────────────────────────────────────────────────
\section{The Falsification Criterion}
\label{sec:falsification}

The falsification criterion is independent of the signal magnitude prediction and holds
regardless of how $\asc$ and the Planck-scale renormalization are resolved.

\begin{keyresult}
\textbf{Falsification conditions F1--F5.}  The framework is falsified for laboratory-scale
relevance if \emph{all} of the following hold simultaneously:

\begin{enumerate}
  \item[\textbf{F1}] Quantum-coherent systems with $\geq 10^{6}$ genuinely entangled qubits
        are achieved, with entanglement depth $\geq 10^{3}$ verified by witnesses.
  \item[\textbf{F2}] Pressure sensitivity reaches $\Delta p < 10^{-6}$\,Pa.
  \item[\textbf{F3}] At least 1000 independent experimental runs yield statistically
        consistent null results.
  \item[\textbf{F4}] Null results are confirmed on $\geq 3$ independent platforms (e.g.,
        trapped neutral atoms, superconducting qubits, optomechanical oscillators).
  \item[\textbf{F5}] No stress-energy contribution beyond standard decoherence models is
        detected at the F2 sensitivity level.
\end{enumerate}

Consequence: $|\kt| < 10^{-15}$, ruling out laboratory-scale relevance for the foreseeable
future.
\end{keyresult}

\paragraph{Why this criterion is strong.}
It is quantitative, platform-independent, statistically rigorous, achievable with current and
near-term technology, and admits no escape hatches (entanglement depth and sensitivity are
fully specified).  A null result satisfying F1--F5 is not a failure --- it is a precise bound
on a previously unmeasured quantity.

% ─────────────────────────────────────────────────────────────────────────────
\section{Open Problems}
\label{sec:open}

\paragraph{OP3 (Critical): Lindblad derivation of $\asc$.}
Derive $\asc$ from a microscopic open-quantum-system model with explicit Hamiltonian, bath,
and coupling.  Without this, predicted signal sizes are indeterminate and cannot be compared
against any experimental bound.  The macroscopic BEC calculation ($\asc \sim 10^{-43}$) and
the NISQ heuristic ($\asc \sim 10^{-4}$--$10^{-2}$) may reflect regime differences rather
than contradiction; a Lindblad derivation would resolve which regime applies.

\paragraph{OP8 (High): Bianchi identity compatibility.}
The contracted Bianchi identity $\nabla^{\mu}G_{\mu\nu} = 0$ implies
$\nabla^{\mu}T_{\mu\nu}^{\mathrm{matter}} = -(\kt c^{4}/8\pi G\kB\lnTwo)\nabla_{\nu}\Sent$.
Standard matter is therefore not independently energy-momentum conserved when $\nabla\Sent
\neq 0$.  Estimate the magnitude of this anomalous force for the proposed NISQ experiment and
compare against precision tests of energy-momentum conservation.

\paragraph{OP10 (Critical, long-term): Planck-scale renormalization.}
The coupling prefactor $c^{4}/(8\pi\kB\lnTwo) \approx 5\times10^{65}$ in SI units reflects
the Planck-scale origin of the derivation.  Whether $\kt = -1/4$ at the Planck scale
renormalizes to a far smaller effective value at laboratory scales --- and what the
renormalized value is --- is an open problem in quantum gravity phenomenology.

% ─────────────────────────────────────────────────────────────────────────────
\section{Discussion and Conclusion}
\label{sec:conclusion}

\subsection{Relationship to Established Physics}

This framework is conservative in ontological commitments: classical GR unchanged, standard
quantum mechanics unchanged, no new particles or fields.  The only novel claim is that
macroscopic quantum entanglement sources stress-energy in a way not captured by the standard
$T_{\mu\nu}$.

\subsection{Relationship to Other Programs}

\emph{vs.\ ER$=$EPR:} Maldacena and Susskind~\cite{Maldacena2013} propose entanglement is
equivalent to geometric connectivity (wormholes).  The present framework proposes entanglement
\emph{sources} stress-energy.  These are complementary, not competing.

\emph{vs.\ Verlinde emergent gravity:} Verlinde~\cite{Verlinde2017} proposes gravity as an
entropic force at galactic scales.  The present framework proposes a laboratory-scale
realization with a specific experimental test.

\subsection{Conclusion}

We have constructed a dimensionally consistent framework in which quantum entanglement entropy
density sources spacetime curvature.  The coupling constant $\kt = -1/4$ is estimated under a
Planck-scale holographic regulator hypothesis that is motivated but unproven.  The signal
magnitude is indeterminate pending resolution of the screening factor.  The falsification
criterion is precise, platform-independent, and achievable with near-term technology.

The framework makes no claim that a signal is detectable.  It claims only: either an
entanglement-dependent gravitational effect will be observed under conditions F1--F5, or the
framework is ruled out for laboratory-scale relevance.  Both outcomes constitute scientific
progress.

% ─────────────────────────────────────────────────────────────────────────────
\section*{Acknowledgments}
The author acknowledges AI research partners for assistance with literature synthesis,
dimensional consistency verification, and critical review.  All theoretical content,
experimental design, and scientific judgments are the author's own.

% ─────────────────────────────────────────────────────────────────────────────
\bibliographystyle{unsrtnat}
\begin{thebibliography}{99}

\bibitem{Jacobson1995}
T.~Jacobson,
``Thermodynamics of spacetime: The Einstein equation of state,''
\textit{Phys.\ Rev.\ Lett.}\ \textbf{75}, 1260 (1995).
\href{https://doi.org/10.1103/PhysRevLett.75.1260}{doi:10.1103/PhysRevLett.75.1260}

\bibitem{Verlinde2017}
E.~Verlinde,
``Emergent gravity and the dark universe,''
\textit{SciPost Phys.}\ \textbf{2}, 016 (2017).
\href{https://doi.org/10.21468/SciPostPhys.2.3.016}{doi:10.21468/SciPostPhys.2.3.016}

\bibitem{Bose2017}
S.~Bose, A.~Mazumdar, G.~W.~Morley, H.~Ulbricht, M.~Toro\v{s},
M.~Paternostro, A.~A.~Geraci, P.~F.~Barker, M.~S.~Kim, and G.~Milburn,
``Spin entanglement witness for quantum gravity,''
\textit{Phys.\ Rev.\ Lett.}\ \textbf{119}, 240401 (2017).
\href{https://doi.org/10.1103/PhysRevLett.119.240401}{doi:10.1103/PhysRevLett.119.240401}
\newline
\textit{Note: Author should verify whether the experimental realization appeared in Nature
623 (2023) and add as a separate citation if so.}

\bibitem{Asenbaum2020}
P.~Asenbaum, C.~Overstreet, T.~Kim, J.~Curti, and M.~A.~Kasevich,
``Atom-interferometric test of the equivalence principle at the $10^{-12}$ level,''
\textit{Phys.\ Rev.\ Lett.}\ \textbf{125}, 191101 (2020).
\href{https://doi.org/10.1103/PhysRevLett.125.191101}{doi:10.1103/PhysRevLett.125.191101}

\bibitem{MICROSCOPE2022}
P.~T.~Touboul \textit{et al.}\ (MICROSCOPE Collaboration),
``MICROSCOPE mission: Final results of the test of the equivalence principle,''
\textit{Phys.\ Rev.\ Lett.}\ \textbf{129}, 121102 (2022).
\href{https://doi.org/10.1103/PhysRevLett.129.121102}{doi:10.1103/PhysRevLett.129.121102}

\bibitem{Maldacena2013}
J.~Maldacena and L.~Susskind,
``Cool horizons for entangled black holes,''
\textit{Fortschritte der Physik}\ \textbf{61}, 781 (2013).
\href{https://doi.org/10.1002/prop.201300020}{doi:10.1002/prop.201300020}

\bibitem{MonetteB2026}
K.~Monette,
``Constraint manifolds and the limits of quantum observability,''
Independent Research Program, March 2026.

\bibitem{MonetteC2026}
K.~Monette,
``Landauer's principle in the emergent thermodynamic information framework,''
Independent Research Program, March 2026.

\end{thebibliography}

\end{document}
